%
% mcp-proxy.tex
%
% Z Specification for mcp-proxy
% Formal model of a bidirectional stdio-to-WebSocket MCP proxy
%

\documentclass[a4paper,10pt,fleqn]{article}
\usepackage[margin=1in]{geometry}
\usepackage{fuzz}
\usepackage[colorlinks=true,linkcolor=blue,citecolor=blue,urlcolor=blue]{hyperref}

\begin{document}

\title{mcp-proxy: A Z Specification}
\author{Formal Model of a Bidirectional MCP Proxy}
\date{March 2026}
\hypersetup{
  pdftitle={mcp-proxy: A Z Specification},
  pdfauthor={Punt Labs},
  pdfsubject={Z Specification},
  pdfcreator={fuzz/probcli}
}
\maketitle

\tableofcontents
\newpage

% ───────────────────────────────────────────────────────────────────
\section{Introduction}
% ───────────────────────────────────────────────────────────────────

\texttt{mcp-proxy} bridges Claude Code's MCP stdio transport to a shared
daemon process over WebSocket.  It forwards JSON-RPC messages opaquely---no
parsing, no routing, no transformation.  Two concurrent channels carry
messages in each direction (stdin$\to$daemon, daemon$\to$stdout), with
coordinated shutdown on stdin EOF or daemon disconnect.

This specification models three aspects of the proxy:

\begin{enumerate}
  \item The \textbf{proxy lifecycle}: startup, connection, bridging, and
    shutdown phases.
  \item The \textbf{bridge}: bidirectional message forwarding with ordering
    guarantees and shutdown coordination.
  \item The \textbf{session key resolver}: process tree walking to find the
    topmost \texttt{claude} ancestor.
\end{enumerate}

% ───────────────────────────────────────────────────────────────────
\section{Basic Types}
% ───────────────────────────────────────────────────────────────────

\begin{zed}
  [MSG, URL]
\end{zed}

$MSG$ represents an opaque MCP JSON-RPC message (one line of NDJSON).
$URL$ represents a daemon WebSocket URL.

% ───────────────────────────────────────────────────────────────────
\section{Free Types}
% ───────────────────────────────────────────────────────────────────

\begin{zed}
  ZBOOL ::= ztrue | zfalse
\end{zed}

\begin{zed}
  ProxyPhase ::= %%
  phaseInit | phaseConnecting | phaseBridging | phaseShutdown
\end{zed}

The proxy moves through four phases: initial argument parsing, dialing
the daemon, bidirectional forwarding, and shutdown.

\begin{zed}
  DialResult ::= dialOk | dialRefused | dialTimeout | dialBadURL
\end{zed}

\begin{zed}
  ShutdownCause ::= %%
  causeStdinEOF | causeDaemonGone | causeSignal | causeError
\end{zed}

\begin{zed}
  ChannelStatus ::= chanActive | chanDone | chanError
\end{zed}

Each of the two bridge channels (stdin-to-daemon and daemon-to-stdout)
is independently active, done, or errored.

% ───────────────────────────────────────────────────────────────────
\section{Global Constants}
% ───────────────────────────────────────────────────────────────────

\begin{axdef}
  maxMessageSize : \nat \\
  maxTreeDepth : \nat \\
  dialTimeoutMs : \nat
  \where
  maxMessageSize = 1048576 \\
  maxTreeDepth = 10 \\
  dialTimeoutMs = 5000
\end{axdef}

$maxMessageSize$ is 1\,MB---the scanner buffer and WebSocket read limit.
$maxTreeDepth$ bounds the process tree walk (safety against cycles).
$dialTimeoutMs$ is the WebSocket connection timeout.

% ───────────────────────────────────────────────────────────────────
\section{Proxy Lifecycle}
% ───────────────────────────────────────────────────────────────────

\begin{schema}{Proxy}
  phase : ProxyPhase \\
  daemonURL : URL \\
  sessionKey : \nat_1 \\
  exitCode : \nat
  \where
  exitCode \leq 2 \\
  sessionKey \leq 5 \\
  phase = phaseInit \implies exitCode = 0 \\
  phase = phaseBridging \implies exitCode = 0
\end{schema}

The proxy holds the daemon URL, the resolved session key (always a
positive PID), and an exit code.  While initializing or bridging, exit
code is always~0.

\begin{schema}{InitProxy}
  Proxy'
  \where
  phase' = phaseConnecting \\
  sessionKey' = 1 \\
  exitCode' = 0
\end{schema}

The session key is set to a placeholder; the daemon URL is left
unconstrained (chosen by the environment).  The session key is
resolved before dialing (modeled separately in the session key
section).

% ───────────────────────────────────────────────────────────────────
\subsection{Dial}

\begin{schema}{DialSuccess}
  \Delta Proxy \\
  \Xi Bridge \\
  key? : \nat_1
  \where
  phase = phaseConnecting \\
  key? \leq 5 \\
  phase' = phaseBridging \\
  sessionKey' = key? \\
  daemonURL' = daemonURL \\
  exitCode' = 0
\end{schema}

\begin{schema}{DialFailure}
  \Delta Proxy \\
  \Xi Bridge \\
  result? : DialResult
  \where
  phase = phaseConnecting \\
  result? \neq dialOk \\
  phase' = phaseShutdown \\
  daemonURL' = daemonURL \\
  sessionKey' = sessionKey \\
  exitCode' = 1
\end{schema}

% ───────────────────────────────────────────────────────────────────
\subsection{Shutdown}

\begin{schema}{Shutdown}
  \Delta Proxy \\
  \Xi Bridge \\
  cause? : ShutdownCause
  \where
  phase = phaseBridging \\
  phase' = phaseShutdown \\
  daemonURL' = daemonURL \\
  sessionKey' = sessionKey \\
  cause? = causeStdinEOF \implies exitCode' = 0 \\
  cause? \neq causeStdinEOF \implies exitCode' = 1
\end{schema}

Only stdin EOF produces exit code~0.  Daemon disconnect, signal, or
error all produce exit code~1.

% ───────────────────────────────────────────────────────────────────
\section{Bridge}
% ───────────────────────────────────────────────────────────────────

The bridge models the two concurrent channels as ordered sequences of
messages, with independent status tracking.

\begin{schema}{Bridge}
  stdinQueue : \seq MSG \\
  daemonQueue : \seq MSG \\
  stdinStatus : ChannelStatus \\
  daemonStatus : ChannelStatus \\
  forwarded : \nat \\
  received : \nat
  \where
  forwarded = \# stdinQueue \\
  received = \# daemonQueue \\
  stdinStatus = chanDone \implies daemonStatus \neq chanActive \\
  forwarded \leq 3 \\
  received \leq 3
\end{schema}

$stdinQueue$ records messages forwarded from stdin to the daemon (in
order). $daemonQueue$ records messages received from the daemon and
written to stdout (in order).  The counts are maintained as schema
fields to make the relationship explicit.

The key invariant: once stdin is done, the daemon channel cannot remain
active---the bridge initiates shutdown when stdin reaches EOF.

\begin{schema}{InitBridge}
  Bridge'
  \where
  stdinQueue' = \langle\rangle \\
  daemonQueue' = \langle\rangle \\
  stdinStatus' = chanActive \\
  daemonStatus' = chanActive \\
  forwarded' = 0 \\
  received' = 0
\end{schema}

% ───────────────────────────────────────────────────────────────────
\subsection{Forward Message (stdin $\to$ daemon)}

\begin{schema}{ForwardMessage}
  \Xi Proxy \\
  \Delta Bridge \\
  msg? : MSG
  \where
  phase = phaseBridging \\
  stdinStatus = chanActive \\
  daemonStatus = chanActive \\
  stdinQueue' = stdinQueue \cat \langle msg? \rangle \\
  forwarded' = forwarded + 1 \\
  daemonQueue' = daemonQueue \\
  daemonStatus' = daemonStatus \\
  received' = received \\
  stdinStatus' = stdinStatus
\end{schema}

A message from stdin is appended to the forwarded sequence.
Both channels must be active.

% ───────────────────────────────────────────────────────────────────
\subsection{Receive Message (daemon $\to$ stdout)}

\begin{schema}{ReceiveMessage}
  \Xi Proxy \\
  \Delta Bridge \\
  msg? : MSG
  \where
  phase = phaseBridging \\
  daemonStatus = chanActive \\
  daemonQueue' = daemonQueue \cat \langle msg? \rangle \\
  received' = received + 1 \\
  stdinQueue' = stdinQueue \\
  stdinStatus' = stdinStatus \\
  forwarded' = forwarded
\end{schema}

A message from the daemon is appended to the received sequence.
This operation does not require stdin to be active---push messages
can arrive even after stdin EOF (during the drain window).

% ───────────────────────────────────────────────────────────────────
\subsection{Stdin EOF}

\begin{schema}{StdinEOF}
  \Xi Proxy \\
  \Delta Bridge
  \where
  phase = phaseBridging \\
  stdinStatus = chanActive \\
  stdinStatus' = chanDone \\
  daemonStatus' = chanDone \\
  stdinQueue' = stdinQueue \\
  daemonQueue' = daemonQueue \\
  forwarded' = forwarded \\
  received' = received
\end{schema}

When stdin reaches EOF, both channels transition out of active.
This models the context cancellation that propagates from the scanner
goroutine to the reader goroutine.

% ───────────────────────────────────────────────────────────────────
\subsection{Daemon Disconnect}

\begin{schema}{DaemonDisconnect}
  \Xi Proxy \\
  \Delta Bridge
  \where
  phase = phaseBridging \\
  daemonStatus = chanActive \\
  daemonStatus' = chanError \\
  stdinStatus' = chanDone \\
  stdinQueue' = stdinQueue \\
  daemonQueue' = daemonQueue \\
  forwarded' = forwarded \\
  received' = received
\end{schema}

When the daemon disconnects unexpectedly, the daemon channel errors
and stdin is closed (the reader goroutine closes stdin to unblock
the scanner).

% ───────────────────────────────────────────────────────────────────
\section{Session Key Resolution}
% ───────────────────────────────────────────────────────────────────

The session key resolver walks a process tree to find the topmost
\texttt{claude} ancestor.  This is a pure function invoked once at
startup---it is not part of the runtime state machine.  We document
its contract here for completeness.

\textbf{Contract.}  Given a process table (mapping PIDs to parent
PIDs) and a set of PIDs whose \texttt{comm} basename is
\texttt{claude}, the resolver returns the topmost claude ancestor
of the proxy's PID.  If no claude ancestor exists, it returns
\texttt{os.Getppid()} as a fallback.  The walk is bounded to
$maxTreeDepth$ iterations (safety against cycles).  On macOS,
\texttt{path.Base(comm)} normalises full paths
(\texttt{/Applications/Claude.app/.../claude}) to basenames.

% ───────────────────────────────────────────────────────────────────
\section{Unified State}
% ───────────────────────────────────────────────────────────────────

\begin{schema}{State}
  Proxy \\
  Bridge
  \where
  phase = phaseBridging \implies %%
  (stdinStatus = chanActive \lor daemonStatus = chanActive) \\
  phase = phaseShutdown \implies %%
  stdinStatus \neq chanActive
\end{schema}

While bridging, at least one channel is active.  After shutdown,
stdin is never active.

\begin{schema}{Init}
  State'
  \where
  phase' = phaseConnecting \\
  sessionKey' = 1 \\
  exitCode' = 0 \\
  stdinQueue' = \langle\rangle \\
  daemonQueue' = \langle\rangle \\
  stdinStatus' = chanActive \\
  daemonStatus' = chanActive \\
  forwarded' = 0 \\
  received' = 0
\end{schema}

% ───────────────────────────────────────────────────────────────────
\section{System Invariants}
% ───────────────────────────────────────────────────────────────────

\begin{enumerate}
  \item \textbf{Message ordering.}  Messages are forwarded in the order
    they arrive on stdin.  Messages from the daemon are written to stdout
    in the order they arrive on the WebSocket.  The sequences
    $stdinQueue$ and $daemonQueue$ capture this.

  \item \textbf{Message opacity.}  The proxy never parses, validates, or
    transforms messages.  Every $MSG$ forwarded is byte-identical to the
    $MSG$ received.

  \item \textbf{Coordinated shutdown.}  Once stdin reaches EOF, the daemon
    channel exits active.  Once the daemon disconnects, stdin is closed.
    Neither channel can remain active alone indefinitely.

  \item \textbf{Clean exit on EOF.}  Exit code~0 only when stdin EOF
    caused shutdown (the caller closed the pipe).

  \item \textbf{Session identity.}  The session key is a positive integer
    (PID) resolved once at startup and passed to the daemon.  It does not
    change during the proxy's lifetime.

  \item \textbf{Bounded tree walk.}  The session key resolver terminates
    in at most $maxTreeDepth$ iterations regardless of process table
    structure (cycles, deep chains).
\end{enumerate}

\end{document}
